\chapter{Նախագծի համատեքստ}

Սույն առաջադրանքը սահմանում է \textbf{միջամբիոնային բակալավրիատի դիպլոմային
նախագիծ}, որը կատարվում է \textbf{երկու ուսանողի} կողմից համատեղ.

\begin{itemize}[leftmargin=1.4em]
  \item \textbf{Ուսանող 1 (Միկրոէլեկտրոնիկա)} — թվային նախագծման / FPGA բաղադրիչ։
  \item \textbf{Ուսանող 2 (Ծրագրային / Համակարգչային ճարտարագիտություն,
        Ավտոմատացում, Ալգորիթմներ կամ հարակից ամբիոն)} — ծրագրային ապահովման,
        տվյալների վերլուծության, վիզուալիզացիայի և մոդելավորման բաղադրիչ։
\end{itemize}

Ապարատային հենքը՝
\begin{itemize}[leftmargin=1.4em]
  \item FPGA հարթակ՝ \textbf{AMD AUP ZU3},
  \item սարքերի ընտանիք՝ \textbf{AMD Zynq UltraScale+ MPSoC},
  \item նախագծման գործիքներ՝ \textbf{AMD Vivado} և \textbf{Vitis}։
\end{itemize}

Նախագիծը պետք է պարունակի էական թվային նախագծման / FPGA բաղադրիչ առաջին ուսանողի
համար և միևնույն ժամանակ ապահովի \textbf{անկախ} ծրագրային, տվյալների վերլուծության,
վիզուալիզացիայի ու մոդելավորման խնդիր երկրորդ ուսանողի համար։ Համակարգը պետք է
աջակցի \textbf{երկկողմանի հաղորդակցություն}.

\begin{enumerate}[leftmargin=1.6em]
  \item \textbf{FPGA $\rightarrow$ Հարթակ.} հավաքագրման և մշակման արդյունքների փոխանցում,
  \item \textbf{Հարթակ $\rightarrow$ FPGA.} աշխատանքային ռեժիմում (runtime) կոնֆիգուրացիա և կառավարում։
\end{enumerate}
